\subsection{ピタゴラス数}

$t=\frac nm$を整数$m,n$で表すと,
\[
X=\frac{m^2-n^2}{m^2+n^2},\qquad
Y=\frac{2mn}{m^2+n^2}.
\]
従って
\begin{equation}\label{式:ピタゴラス数}
(m^2-n^2)^2+(2mn)^2=(m^2+n^2)^2
\end{equation}
である.\,$m>n>0$なら三つの整数は正であり,直角三角形の三辺になる.\,$a^2+b^2=c^2$を満たす正の整数の組を\textbf{ピタゴラス数}という.

逆に整数の直角三角形$(a,b,c)$があれば,点$\left(\frac ac,\frac bc\right)$は単位円上にある.その傾きは
\[
t=\cfrac{\cfrac bc}{1+\cfrac ac}=\frac{b}{c+a}
\]
と有理数になる.これを$\frac nm$と表せば,上の公式が元の三辺と同じ比を与える.従って,この表示は偶然いくつかの例を作るだけでなく,すべての整数の直角三角形の辺の比を表す.得られた三整数に共通因数があれば,それで割ればよい.

三辺の最大公約数が$1$のピタゴラス数を,\textbf{原始ピタゴラス数}という.整数$m>n>0$が互いに素で,一方が偶数,他方が奇数ならば,式\eqref{式:ピタゴラス数}は原始ピタゴラス数を与える.逆に,すべての原始ピタゴラス数は,直角を挟む二辺の交換を除いてこの形で表せる.この主張を,整数の素因数分解の一意性を使って確かめよう.

まず上の条件で作った三整数に,共通の素因数$p$があると仮定する.\,$m^2-n^2$と$m^2+n^2$はともに奇数なので,\,$p$は奇素数である.この二数の和と差から,\,$p$は$2m^2$と$2n^2$を割る.従って$p$は$m,n$の両方を割り,互いに素という条件に反する.共通の素因数がないので,三整数の最大公約数は$1$である.

逆に原始ピタゴラス数$(a,b,c)$を考える.二辺$a,b$がともに偶数なら原始ではなく,ともに奇数なら$a^2+b^2$を$4$で割った余りが$2$となり,平方数$c^2$にはなれない.従って一方だけが偶数であり,必要なら辺を交換して$b$を偶数,\,$a$を奇数とする.このとき$c$も奇数である.また$a,c$に共通の素因数があれば,\,$b^2=c^2-a^2$,従って$b$もその素因数で割れ,原始という条件に反する.よって$a,c$は互いに素である.

ここで
\[
A=\frac{c+a}{2},\qquad B=\frac{c-a}{2}
\]
と置くと,\,$A>B>0$の整数であり,
\[
AB=\frac{c^2-a^2}{4}=\left(\frac b2\right)^2.
\]
$A,B$の共通因数は$A+B=c$と$A-B=a$も割るので,\,$A,B$は互いに素である.互いに素な二数の積が平方数なら,素因数は片方にしか現れず,その指数はすべて偶数になる.従って$A=m^2$,\,$B=n^2$と書ける.和と差を戻せば
\[
a=m^2-n^2,\qquad b=2mn,\qquad c=m^2+n^2.
\]
$A,B$が互いに素だから$m,n$も互いに素であり,\,$c$が奇数だから$m,n$の偶奇は異なる.これで逆向きも示せた.円の表示が辺の比を与えることに加えて,約分の要らない整数を作る条件まで分かった.

\begin{example}{ピタゴラス数}
$t=\frac23$からピタゴラス数を作れ.また三辺が$8,15,17$の三角形に対応する傾きを求めよ.
\end{example}

$m=3,n=2$なので
\[
m^2-n^2=5,\qquad2mn=12,\qquad m^2+n^2=13.
\]
確かに$5^2+12^2=13^2$である.後半では$P=\left(\frac{15}{17},\frac8{17}\right)$と選ぶと
\[
t=\frac8{17+15}=\frac14.
\]
$m=4,n=1$を公式へ戻せば$(15,8,17)$となる.二つの直角を挟む辺のどちらを$X$座標にするかで,対応する傾きは変わる.



